%=================================================================
\vspace{6pt}

% TODO(author): confirm every statement below before submission.
\authorcontributions{Conceptualization, H.S. and B.D.; methodology, H.S.; software, H.S.; validation, H.S., T.L. and W.S.; formal analysis, H.S.; investigation, T.L. and J.L. (Jie Lin); resources, B.D.; data curation, W.S.; writing---original draft preparation, H.S.; writing---review and editing, B.D. and J.L. (Junsong Luo); visualization, H.S.; supervision, B.D.; project administration, B.D.; funding acquisition, B.D. All authors have read and agreed to the published version of the~manuscript.}

\funding{This research received no external~funding.}

\institutionalreview{Not applicable.}

\informedconsent{Not applicable.}

% TODO(author): 采用的是建议措辞，投稿前请确认是否符合实际数据开放情况。
\dataavailability{The simulation code and the data generated in this study are available from the corresponding author upon reasonable~request.}

\conflictsofinterest{The authors declare no conflicts of~interest.}

\abbreviations{Abbreviations}{%
The following abbreviations are used in this manuscript:\\

\noindent
\begin{tabular}{@{}ll}
5G & Fifth generation\\
ADH & Adaptive dynamic hysteresis\\
ALERA & Attention-LSTM enhanced risk-aware vertical handoff algorithm\\
COPRAS & Complex proportional assessment\\
GRA & Grey relational analysis\\
gNB & Next generation Node B\\
LAAVHA & LSTM-attention based adaptive vertical handover algorithm\\
LCB & Lower confidence bound\\
LSTM & Long short-term memory\\
LTE & Long-Term Evolution\\
MADM & Multiple attribute decision making\\
PLR & Packet loss ratio\\
RSRP & Reference signal received power\\
RS-TOPSIS & Risk-sensitive TOPSIS\\
SAW & Simple additive weighting\\
SINR & Signal-to-interference-plus-noise ratio\\
SPOTIS & Stable preference ordering toward ideal solution\\
TOPSIS & Technique for order preference by similarity to the ideal solution\\
UAV & Unmanned aerial vehicle\\
VHO & Vertical handover\\
VIKOR & VIsekriterijumska Optimizacija I Kompromisno Resenje
\end{tabular}
}

%=================================================================
\begin{adjustwidth}{-\extralength}{0cm}

\reftitle{References}

\begin{thebibliography}{999}
\bibitem{ref1}
Yang, K.; Wang, Y.; Gao, X.; et al. Communications in space--air--ground integrated networks: An overview. {\em Space Sci. Technol.} {\bf 2025}, {\em 5}, 0199.
\bibitem{ref2}
Geraci, G.; García-Rodríguez, A.; Azari, M.M.; et al. What will the future of UAV cellular communications be? A flight from 5G to 6G. {\em IEEE Commun. Surv. Tutor.} {\bf 2022}, {\em 24}, 1304--1335.
\bibitem{ref3}
Ribeiro, L.M.B.; Müller, I.; Becker, L.B. Communication interface manager for improving performance of heterogeneous UAV networks. {\em Sensors} {\bf 2021}, {\em 21}, 4255.
\bibitem{ref4}
Ayass, T.; Coqueiro, T.; Carvalho, T.; et al. Unmanned aerial vehicle with handover management fuzzy system for 5G networks: Challenges and perspectives. {\em Intell. Robot.} {\bf 2022}, {\em 2}, 20--36.
\bibitem{ref5}
Rehman, A.U.; Roslee, M.B.; Tiang, J.J. A survey of handover management in mobile HetNets: Current challenges and future directions. {\em Appl. Sci.} {\bf 2023}, {\em 13}, 3367.
\bibitem{ref6}
Haghrah, A.; Abdollahi, M.P.; Azarhava, H.; et al. A survey on the handover management in 5G-NR cellular networks: Aspects, approaches and challenges. {\em EURASIP J. Wirel. Commun. Netw.} {\bf 2023}, {\em 2023}, 52.
\bibitem{ref7}
Wang, Z.; Lv, Z.; Xu, X.; et al. Vertical switching algorithm for unmanned aerial vehicle in power grid heterogeneous communication networks. {\em Electronics} {\bf 2024}, {\em 13}, 2612.
\bibitem{ref8}
Zhong, J.; Zhang, L.; Alhabo, M.; et al. A hybrid scheme using TOPSIS and Q-learning for handover decision making in UAV-assisted heterogeneous network. {\em IEEE Access} {\bf 2024}, {\em 12}, 31422--31430.
\bibitem{ref9}
Goh, M.I.; Mbulwa, A.I.; Yew, H.T.; et al. Handover decision-making algorithm for 5G heterogeneous networks. {\em Electronics} {\bf 2023}, {\em 12}, 2384.
\bibitem{ref10}
Hwang, C.L.; Yoon, K. \textit{Multiple Attribute Decision Making: Methods and Applications. A State-of-the-Art Survey}; Springer-Verlag: Berlin, Germany, 1981.
\bibitem{ref11}
Zolfani, S.; Yazdani, M.; Pamucar, D.; et al. A VIKOR and TOPSIS focused reanalysis of the MADM methods based on logarithmic normalization. {\em Facta Univ. Ser. Mech. Eng.} {\bf 2020}, {\em 18}, 341--355.
\bibitem{ref12}
Dantas Silva, F.S.; Lima, M.P.S.; Corujo, D.; et al. A comprehensive step-wise survey of multiple attribute decision-making mobility approaches. {\em IEEE Access} {\bf 2024}, {\em 12}, 108616--108656.
\bibitem{ref13}
Satapathy, P.; Mahapatro, J. An efficient multicriteria-based vertical handover decision-making algorithm for heterogeneous networks. {\em Trans. Emerg. Telecommun. Technol.} {\bf 2022}, {\em 33}, e4409.
\bibitem{ref14}
Zhong, Y.; Wang, H.; Lv, H.; et al. A vertical handoff decision scheme using subjective-objective weighting and grey relational analysis in cognitive heterogeneous networks. {\em Ad Hoc Netw.} {\bf 2022}, {\em 134}, 102924.
\bibitem{ref15}
Riaz, H.; Öztürk, S.; Aldirmaz Çolak, S.; et al. Performance analysis of weighting methods for handover decision in HetNets. {\em Gazi Univ. J. Sci.} {\bf 2024}, {\em 37}, 1791--1810.
\bibitem{ref16}
Tan, K.; Bremner, D.; Le Kernec, J.; et al. Intelligent handover algorithm for vehicle-to-network communications with double-deep Q-learning. {\em IEEE Trans. Veh. Technol.} {\bf 2022}, {\em 71}, 7848--7862.
\bibitem{ref17}
Jang, Y.; Raza, S.M.; Kim, M.; et al. Proactive handover decision for UAVs with deep reinforcement learning. {\em Sensors} {\bf 2022}, {\em 22}, 1200.
\bibitem{ref18}
Zaid, M.; Kadir, M.K.A.; Shayea, I.; et al. Machine learning-based approaches for handover decision of cellular-connected drones in future networks: A comprehensive review. {\em Eng. Sci. Technol. Int. J.} {\bf 2024}, {\em 55}, 101732.
\bibitem{ref19}
Hochreiter, S.; Schmidhuber, J. Long short-term memory. {\em Neural Comput.} {\bf 1997}, {\em 9}, 1735--1780.
\bibitem{ref20}
Kaur, G.; Goyal, R.K.; Mehta, R. An efficient handover mechanism for 5G networks using hybridization of LSTM and SVM. {\em Multimed. Tools Appl.} {\bf 2022}, {\em 81}, 37057--37085.
\bibitem{ref21}
Alraih, S.; Nordin, R.; Abu-Samah, A.; et al. A survey on handover optimization in beyond 5G mobile networks: Challenges and solutions. {\em IEEE Access} {\bf 2023}, {\em 11}, 59317--59345.
\bibitem{ref22}
Fan, D.; Zhou, S.; Luo, J.; et al. Joint traffic prediction and handover design for LEO satellite networks with LSTM and attention-enhanced rainbow DQN. {\em Electronics} {\bf 2025}, {\em 14}, 3040.
\bibitem{ref23}
Zhou, S.; Liu, X.; Tang, B.; et al. Handover and coverage analysis in 3-D mobile UAV cellular networks. {\em IEEE Internet Things J.} {\bf 2024}, {\em 11}, 29911--29925.
\bibitem{ref24}
Tashan, W.; Shayea, I.; Aldirmaz-Colak, S.; et al. Voronoi-based handover self-optimization technique for handover ping-pong reduction in 5G networks. In Proceedings of the 2023 10th International Conference on Wireless Networks and Mobile Communications (WINCOM), Istanbul, Türkiye, 26--28 October 2023; pp. 1--6.
\bibitem{ref25}
Ndegwa, S.; Nyachionjeka, K.; Mharakurwa, E.T. User preference-based heterogeneous network management system for vertical handover. {\em J. Electr. Comput. Eng.} {\bf 2023}, {\em 2023}, 1--20.
\bibitem{ref26}
Soydaner, D. Attention mechanism in neural networks: Where it comes and where it goes. {\em Neural Comput. Appl.} {\bf 2022}, {\em 34}, 13371--13385.
\bibitem{ref27}
Vaswani, A.; Shazeer, N.; Parmar, N.; et al. Attention is all you need. In Proceedings of the Advances in Neural Information Processing Systems 30 (NIPS 2017), Long Beach, CA, USA, 4--9 December 2017; pp. 5998--6008.
\bibitem{ref28}
Hwang, W.S.; Cheng, T.Y.; Wu, Y.J.; et al. Adaptive handover decision using fuzzy logic for 5G ultra-dense networks. {\em Electronics} {\bf 2022}, {\em 11}, 3278.
\bibitem{ref29}
Gannapathy, V.R.; Nordin, R.; Abu-Samah, A.; et al. An adaptive TTT handover (ATH) mechanism for dual connectivity (5G mmWave--LTE advanced) during unpredictable wireless channel behavior. {\em Sensors} {\bf 2023}, {\em 23}, 4357.
\bibitem{ref30}
Karmakar, R.; Kaddoum, G.; Chattopadhyay, S. Mobility management in 5G and beyond: A novel smart handover with adaptive time-to-trigger and hysteresis margin. {\em IEEE Trans. Mob. Comput.} {\bf 2023}, {\em 22}, 5995--6010.
\bibitem{ref31}
Ezz-Eldien, N.A.; Abdel-Atty, H.M.; Abdalla, M.I.; et al. An adaptive optimized handover decision model for heterogeneous networks. {\em PLoS ONE} {\bf 2023}, {\em 18}, e0294411.
\bibitem{ref32}
Manzoor, S.; Manzoor, M.; Manzoor, H.; et al. Which simulator to choose for next generation wireless network simulations? ns-3 or OMNeT++. {\em Eng. Proc.} {\bf 2023}, {\em 46}, 36.
\bibitem{ref33}
Yin, H.; Liu, P.; Liu, K.; et al. ns3-ai: Fostering artificial intelligence algorithms for networking research. In Proceedings of the 2020 Workshop on ns-3 (WNS3 '20), Gaithersburg, MD, USA, 17--18 June 2020; ACM: New York, NY, USA, 2020; pp. 57--64.
\bibitem{ref34}
Rappaport, T.S. \textit{Wireless Communications: Principles and Practice}, 2nd ed.; Prentice Hall: Upper Saddle River, NJ, USA, 2002.
\end{thebibliography}

\PublishersNote{}
\end{adjustwidth}
\end{document}
